% ============================================================================
\clearpage
\section*{v4.9.0-r33 Source-Only: Left-Invariant Labels on the Triple Space and a Pre-Registered Test}
\addcontentsline{toc}{section}{v4.9.0-r33 Source-Only: Left-Invariant Labels on the Triple Space and a Pre-Registered Test}
\label{sec:v490-r33-left-invariant-labels}
% ============================================================================

\paragraph{Purpose.}
Layer r32 closed the route of natural potentials in the slot norms and real parts.  On the author's request the Dixon
algebra \(\mathbb C\otimes\HH\otimes\OO\) was then read as the storage for all configurations that the Companion intended
it to be (the HS(6) carrier of the Dixon-carrier coding layer), and the generation triple was placed in it.  This layer
records what follows for the generation triple, a resolution of the conflict with the Lorentz spinor, and the
pre-registered test H18 (note \texttt{llm/Vorab\_Notiz\_H18\_Links\_invariante\_Massen.md}, commit \texttt{6130814}, fixed
before computing).  The algebraic statements are proved in Lean (\texttt{formal/QSFormal/Dixon.lean} and
\texttt{Lesarten.lean}, Lean~4.35.0, no Mathlib import, standard axioms only).  Earlier layers are not rewritten.

\noindent\textbf{Source-only status.}
This v4.9.0-r33 layer is source-only.  It does not build a PDF, does not assign a new DOI,
and does not authorize a public release.  The v4.9.0 Version DOI
\href{https://doi.org/10.5281/zenodo.23534954}{10.5281/zenodo.23534954} and the Concept DOI
\href{https://doi.org/10.5281/zenodo.19210728}{10.5281/zenodo.19210728} are unchanged by it.

\subsection*{v4.9.0-r33 the triple in the Dixon storage}
\label{subsec:v490-r33-storage}

\paragraph{Theoremlet v4.9.0-r33.1.}
\begin{enumerate}[leftmargin=2em]
\item[(a)] \(\HH\otimes\OO=\OO\oplus(\operatorname{Im}\HH)\otimes\OO\), and \((\operatorname{Im}\HH)\otimes\OO=i\otimes\OO\oplus
j\otimes\OO\oplus k\otimes\OO\) carries the triple of v4.8.0-r14/r15: 24 of the 64 channels, with the coloured states as a
9-channel submask.  \(G_2\) acts on \(\OO\), \(\operatorname{Aut}(\HH)=SO(3)\) on \(\HH\); they commute.
\item[(b)] Of the 24 signed permutations of the slots in \(SO(3)\), exactly the 12 even permutations with an even number of
sign changes leave \(\operatorname{Re}((xy)z)\) invariant; they are the conjugations by the binary tetrahedral group, so
the joint symmetry is \(A_4\).  The transpositions of the triality \(S_3\) need octonion conjugation, an anti-automorphism.
\item[(c)] If the \(\mathbb C\) factor is read as the amplitude plane of v4.8.0-r02, then \(J_B=R_{\pi/2}\circ T\) (v4.9.0-r30.2(b))
exchanges its two halves; the slot norms are invariant under \(T\), \(R_{\pi/2}\) and \(J_B\), and the triality form goes to
\(-i\,\bar t\).  Of the rotations of the \(\mathbb C\) plane only the quarter steps, together with the reflections (the
group \(D_4\)), keep the label of the two halves.
\end{enumerate}

\paragraph{Proof sketch.}
Finite checks and polynomial identities over \(\mathbb Z\), formalised in \texttt{Dixon.lean} (\texttt{inv\_iff},
\texttt{tripleMask\_popcount}, \texttt{jB\_eq\_swap}, \texttt{tC\_JB}) and \texttt{Lesarten.lean} (\texttt{d4\_length},
\texttt{conj\_cyc}, \texttt{conj\_i}).  \(\square\)

\subsection*{v4.9.0-r33 left-invariant labels}
\label{subsec:v490-r33-labels}

\paragraph{Theoremlet v4.9.0-r33.2.}
Write the triple as eight quaternions \(h_c\), one for each octonion component, and let the Lorentz spinor act from the
left, \(h_c\mapsto qh_c\).  The quaternionic Gram matrix \(\overline{h_c}\,h_d\) is invariant (up to \(|q|^2\); with the complex
bilinear form also under boosts).  Its real part is symmetric in the slots.  Its three imaginary components are three
invariant labels on the right.  For the pure triple the \(i\) component is \(-x_j\wedge x_k\), so
\[
\mu_i=2\bigl(|x_j|^2|x_k|^2-\langle x_j,x_k\rangle^2\bigr),
\]
twice the area spanned by the other two slots.  The right action of \(1+i+j+k\) permutes the three labels cyclically.  For
orthogonal slots \(\sqrt{\mu_i/2}=n_jn_k=P/n_i\).

\paragraph{Proof sketch.}
Associativity of \(\HH\) gives \(\overline{qx}\,qy=|q|^2\,\bar x y\) and \(\overline{xr}\,yr=\bar r(\bar xy)r\); for imaginary
\(u,v\), \(\bar uv=(u\cdot v,-u\times v)\); the Lagrange identity in \(\mathbb R^8\) gives \(\mu_i\).  Lean: \texttt{gram\_left},
\texttt{gram\_right}, \texttt{right\_cyc}, \texttt{gram\_pure}, \texttt{lagrange8}; numerical check
\texttt{llm/k72\_links\_invariant.py}.  \(\square\)

\noindent With Lorentz acting from the left the single slot norms are not invariant (a rotation redistributes them), so the
conflict of the Dixon check with the spinor reading of Theorem~\ref{thm:three-generations} can be resolved for the labels by
reading the generations on the right; it is not resolved for masses (see H18).

\subsection*{v4.9.0-r33 pre-registered test H18}
\label{subsec:v490-r33-h18}

H18: masses \(m_k\propto\sqrt{\mu_k}\) at the global minima of the potentials already tested (family A: v4.8.0-r16;
family B: r32), with a fixed parameter grid and distribution.  H18a: some point gives determined, distinct, non-zero
\(\mu_k\); H18b: one of them meets both charged-lepton ratios within \(1\%\).  Script \texttt{llm/k73\_test\_h18.py}: family A,
205 points, all undetermined; family B, 500 points, 428 undetermined, 51 with equal and 21 with two distinct values.
\textbf{H18a and H18b are refuted.}  Check after the verdict (\texttt{llm/k74\_nachkontrolle\_h18.py}): in family A the angle
between two slots is a flat direction of the minimum, because \(V\) sees the angles only through \(t\); in family B the
undetermined cases have \(\mu\approx0\), the slots lying parallel on the real axis.  The labels resolve the Lorentz conflict
but inherit the selection problem: a mass hierarchy would need a principle that fixes the angles between the slots.

\subsection*{v4.9.0-r33 status in the three frames}
\label{subsec:v490-r33-frames}

\begin{center}
\small
\begin{tabular}{@{}>{\raggedright\arraybackslash}p{0.30\textwidth}>{\raggedright\arraybackslash}p{0.62\textwidth}@{}}
\toprule
result & M / T / P\\
\midrule
r33.1(a), (b) & algebraic, shown in the reference frame R in which the Dixon grammar is written (class N of r23)\\
r33.1(c) & under the reading of the \(\mathbb C\) factor as amplitude plane: J\(_B\) exchanges the halves (M \(\leftrightarrow\) T); slot norms frame-free; P has equal weight on both halves\\
r33.2, H18 & frame-free in the slots; the \(\mathbb C\) factor does not enter\\
\bottomrule
\end{tabular}
\end{center}

\subsection*{v4.9.0-r33 register entries}
\label{subsec:v490-r33-register}

\begin{center}
\small
\begin{tabular}{@{}>{\raggedright\arraybackslash}p{0.34\textwidth}>{\raggedright\arraybackslash}p{0.58\textwidth}@{}}
\toprule
statement & status after v4.9.0-r33\\
\midrule
Dixon-carrier coding (HS(6)) & read as the storage for all configurations (author's intention); holds the generation triple; masks conservative\\
mixed readouts (guarded, binary masks) & defined, never evaluated; r33 does not evaluate one either\\
Theorem~\ref{thm:three-generations} & colour and generation on the same triplet (r13); spinor reading in conflict with generations in \(\operatorname{Im}\HH\) unless the generations are read on the right (r33.2)\\
mass hierarchy on the triple space & refuted also through left-invariant labels (H18); needs a principle that fixes the angles between the slots\\
\bottomrule
\end{tabular}
\end{center}

\subsection*{v4.9.0-r33 ledger and guard update}
\label{subsec:v490-r33-ledger-guard}

\[
R_{\rm cg}^{v490\text{-}r33}=(0,0,1)
\]
for (separation consistent with \(G_2\); carrier and storage of the triality slots identified, Lorentz-compatible labels found;
mass mechanism open).

The layer spends no guard:
\[
R_{\rm glob}^{(v490\text{-}r33,+)}=(0,1,1,1,1),
\qquad
\Gamma_{\rm guard}^{v490\text{-}r33}=(1,1,1,1,1),
\qquad
\beta_4=1/256\ \text{unspent}.
\]
It does not imply exact finite \(C5\) zero closure, exact all-readout terminality,
\(C6/T2\), a global photon-clock/RH theorem, or global all-readout closure.
